\documentclass[11pt]{article}
\usepackage[T1]{fontenc}
\usepackage[margin=1in]{geometry}
\usepackage{amsmath,amssymb}
\usepackage{enumitem}
\usepackage{hyperref}
\hypersetup{colorlinks=true,linkcolor=black,urlcolor=blue}

\begin{document}

\begin{center}
{\LARGE HRV Formula Sheet}\\[0.5em]
{\large PolarH10 short-term RR-derived HRV}
\end{center}

\section*{Data path}

\begin{itemize}[leftmargin=1.5em]
\item Input stream: accepted RR intervals from live BLE heart-rate measurement packets.
\item Physical H10: RR intervals are derived inside the strap from ECG beat timing.
\item Synthetic transport: RR remains the source for the HRV solve, while synthetic PMD ECG is emitted on the same synthetic beat schedule for operator review.
\end{itemize}

\section*{Rolling short-term window}

The app uses a rolling RR window with a default length of $300$ s. That default
tracks the conventional short-term setting discussed by Shaffer and Ginsberg.
The first solve is intentionally delayed until the rolling window is nearly full
enough to avoid publishing an early headline value from a partial capture.

\section*{Time-domain formulas}

For accepted RR intervals $NN_1 \ldots NN_n$:

Here $NN_i$ means the $i$-th accepted beat interval in milliseconds, $n$ is the
number of accepted intervals in the current rolling window, and an overline
means a window average. These formulas are applied to the accepted short-term
series the tracker keeps for the HRV solve, not to every raw packet value.

\[
\overline{NN} = \frac{1}{n}\sum_{i=1}^{n} NN_i
\]

This is the average accepted beat interval over the current short-term window.

\[
\overline{HR} = \frac{60000}{\overline{NN}}
\]

This converts the mean interval into beats per minute using $60000$ ms $= 1$
minute.

\[
SDNN = \sqrt{\frac{1}{n-1}\sum_{i=1}^{n}(NN_i - \overline{NN})^2}
\]

$SDNN$ measures how widely the accepted intervals are spread across the whole
window.

\[
RMSSD = \sqrt{\frac{1}{n-1}\sum_{i=2}^{n}(NN_i - NN_{i-1})^2}
\]

$RMSSD$ focuses on beat-to-beat change by differencing neighboring intervals
before averaging.

\[
\ln RMSSD = \ln(RMSSD)
\]

$\ln RMSSD$ is the natural-log form of $RMSSD$; it compresses the scale and is
often easier to compare across sessions.

\[
pNN50 = 100 \times
\frac{\#\left\{|NN_i - NN_{i-1}| > 50\,ms\right\}}{n-1}
\]

$pNN50$ is the percentage of adjacent interval pairs whose absolute difference
exceeds $50$ ms.

\[
SD1 = \frac{RMSSD}{\sqrt{2}}
\]

$SD1$ is the short-axis Poincare equivalent of short-term variability and is
shown because many readers know HRV through that geometric view.

\section*{Alignment status}

The implementation is aligned to the cited short-term HRV guidance for RMSSD,
SDNN, pNN50 with a fixed $50$ ms threshold, SD1, mean NN, mean HR, and
$\ln RMSSD$.

The app does not claim to reproduce 24 h norms, clinical artifact adjudication,
or full Holter-style NN editing. It is therefore documented as short-term
RR-derived telemetry rather than a clinical HRV workstation.

\section*{References}

\begin{itemize}[leftmargin=1.5em]
\item Shaffer and Ginsberg, ``An Overview of Heart Rate Variability Metrics and Norms'' (2017).
\item \url{https://doi.org/10.3389/fpubh.2017.00258}
\item \url{https://mesmerprism.github.io/PolarH10/reference/hrv-workflow.html}
\item \url{https://mesmerprism.github.io/PolarH10/reference/references.html}
\end{itemize}

\end{document}
